\chapterbackmatter{Solutions to Supplementary Exercises}

\begin{enumerate}
  \SupplementarySolution{1}{The Mass Matrix as a Gram Matrix}
  Because \(t_\alpha\) is imaginary and antisymmetric, define the real
  vectors \(w_\alpha=i t_\alpha v\).  Then
  \[
    \mu^2_{\alpha\beta}=w_\alpha\cdot w_\beta.
  \]
  For arbitrary real \(c_\alpha\),
  \[
    c_\alpha\mu^2_{\alpha\beta}c_\beta
      =\left|\sum_\alpha c_\alpha w_\alpha\right|^2\geq0.
  \]
  Thus \(\mu^2\) is a Gram matrix and is positive semidefinite.  Equality
  holds precisely when
  \[
    \left(\sum_\alpha c_\alpha t_\alpha\right)v=0,
  \]
  which says that the corresponding Lie-algebra element stabilizes the
  vacuum.  Therefore
  \[
    \ker\mu^2=\mathfrak h,\qquad
    \operatorname{rank}\mu^2=\dim\mathfrak g-\dim\mathfrak h.
  \]
  Each nonzero eigenvalue is the squared mass of one gauge boson, so the
  rank counts the massive vectors.

  \SupplementarySolution{2}{The Abelian Higgs Spectrum}
  To quadratic order,
  \[
    |D_\mu\Phi|^2
      =\frac12(\partial h)^2+\frac12(\partial\chi)^2
       +\frac12q^2v^2A_\mu A^\mu
       +qv A^\mu\partial_\mu\chi+\cdots
  \]
  and
  \[
    \lambda\left(|\Phi|^2-\frac{v^2}{2}\right)^2
      =\lambda v^2h^2+\cdots.
  \]
  A gauge transformation removes \(\chi\).  The physical masses are
  \[
    \boxed{m_A^2=q^2v^2,\qquad m_h^2=2\lambda v^2.}
  \]
  The \(\chi\) direction is massless in the scalar potential but is not
  a physical particle in unitarity gauge.  Before breaking there are two
  photon polarizations and two real scalar modes, for four degrees of
  freedom.  Afterwards there are three polarizations of the massive
  vector plus the radial scalar, again four.

  \SupplementarySolution{3}{Orbit Directions and Eaten Modes}
  A curve through the orbit is
  \(v(s)=\exp(is^\alpha t_\alpha)v\).  Its tangent at the identity is
  \[
    \left.\frac{dv}{ds^\alpha}\right|_{s=0}=it_\alpha v.
  \]
  Hence the image of the map
  \(\mathfrak g\to V\), \(X\mapsto Xv\), is
  \(T_v(Gv)\).  Its kernel is the stabilizer algebra \(\mathfrak h\), so
  rank--nullity gives
  \[
    \dim T_v(Gv)=\dim G-\dim H.
  \]
  These are exactly the zero modes of the potential Hessian generated by
  broken symmetries.  Gauge transformations remove them, while the same
  vectors \(t_\alpha v\) determine the nonzero gauge-boson mass matrix.
  Thus each removed orbit direction supplies one longitudinal
  polarization.

  \SupplementarySolution{4}{Residual Gauge Freedom in Unitarity Gauge}
  An infinitesimal transformation generated by \(X\) preserves the
  chosen vacuum representative when \(Xv=0\).  These transformations
  form the unbroken algebra \(\mathfrak h\).  Unitarity gauge therefore
  fixes only the broken directions; gauge fields belonging to \(H\)
  remain massless and retain their usual gauge redundancy.  A condition
  such as Lorenz gauge must still be imposed on them to invert their
  quadratic kinetic operator.

  The construction also assumes that the scalar configuration lies on a
  regular orbit with nonzero modulus.  At a zero of the order parameter,
  its angular variables are undefined, so no smooth gauge
  transformation can choose one orbit representative there.  Vortex
  cores are the elementary example.  More generally, distinct gauge
  copies or changing stabilizers can obstruct a single global
  unitary-gauge slice.

  \SupplementarySolution{5}{Poles in an \(R_\xi\) Gauge}
  Adding
  \(-(\partial\cdot A-\xi m\chi)^2/(2\xi)\) cancels the
  \(mA^\mu\partial_\mu\chi\) mixing.  The propagators have the forms
  \[
    \Delta_{\mu\nu}(k)
      =\frac1{k^2+m^2}
       \left(\eta_{\mu\nu}
        \frac{(\xi-1)k_\mu k_\nu}{k^2+\xi m^2}\right),
  \]
  \[
    \Delta_\chi(k)=\frac1{k^2+\xi m^2},
    \qquad
    \Delta_{\rm gh}(k)=\frac1{k^2+\xi m^2}.
  \]
  The pole at \(k^2=-m^2\) is physical.  The common pole at
  \(k^2=-\xi m^2\) is gauge-dependent and cancels from physical
  amplitudes among longitudinal, Goldstone, and ghost contributions.
  Taking \(\xi\to\infty\) removes the last two fields and gives
  \[
    \Delta_{\mu\nu}^{\rm U}(k)
      =\frac{\eta_{\mu\nu}+k_\mu k_\nu/m^2}{k^2+m^2},
  \]
  whose poor large-\(k\) behavior is the familiar unitary-gauge price.

  \SupplementarySolution{6}{A General Tree-Level \(\rho\) Formula}
  For a neutral component, \(T_3=-Y\).  The charged generators obey
  \[
    \frac12\langle T_+T_-+T_-T_+\rangle
      =T(T+1)-T_3^2=T(T+1)-Y^2.
  \]
  This fixes the charged-vector mass.  The neutral mass comes from the
  combination \(gW^3-g'B\) and is proportional to \(2Y^2\).  Summing
  independent vacuum contributions gives
  \[
    \rho_{\rm tree}
      =\frac{\sum_i[T_i(T_i+1)-Y_i^2]|v_i|^2}
             {\sum_i2Y_i^2|v_i|^2}.
  \]
  For \((T,Y)=(\tfrac12,\tfrac12)\), numerator and denominator are both
  \(\tfrac12|v|^2\), so \(\rho_{\rm tree}=1\).

  A \(Y=0\) triplet has a neutral state with \(T_3=0\).  Its vacuum
  contributes to \(m_W\) but not to the neutral mass, so it would make
  the displayed denominator vanish if it were the only multiplet.  For a
  real representation one must also include the conventional factor
  \(1/2\) associated with its field normalization.

  \SupplementarySolution{7}{Why a Chiral Mass Is Forbidden}
  Under the two groups,
  \[
    \bar\psi_L\psi_R\longrightarrow
    e^{\,i[(Q_R-Q_L)\alpha+(Y_R-Y_L)\beta]}
      \bar\psi_L\psi_R.
  \]
  A bare Dirac mass is allowed only if
  \(Q_R=Q_L\) and \(Y_R=Y_L\).  For the stated assignments,
  \[
    Q_R-Q_L=\frac12,\qquad
    Y_R-Y_L=-\frac12,
  \]
  so the bilinear is charged under both independent symmetries.  The
  chiral gauge symmetry forbids a mass until an additional field
  compensates these charges.

  \SupplementarySolution{8}{Solving for Higgs Charges}
  For the interaction \(H\bar\psi_L\psi_R\), invariance requires
  \[
    Q_H-Q_L+Q_R=0,\qquad
    Y_H-Y_L+Y_R=0.
  \]
  The assignments of Exercise S.21.7 therefore give
  \[
    \boxed{(Q_H,Y_H)=(-\tfrac12,+\tfrac12).}
  \]
  If \(H=(v+h)/\sqrt2\) in a gauge where its vacuum value is real, then
  \[
    yH\bar\psi_L\psi_R+\mathrm{H.c.}
      =\frac{yv}{\sqrt2}\bar\psi\psi
       +\frac{y}{\sqrt2}h\bar\psi\psi.
  \]
  Thus the same gauge-invariant coupling generates both the Dirac mass
  \(m=yv/\sqrt2\) and the scalar interaction.

  \SupplementarySolution{9}{An Unbroken Combination of Two \(U(1)\)s}
  The vacuum kinetic term is
  \[
    |D_\mu\langle\Phi\rangle|^2
      =\frac{v^2}{2}
       (g_1q_1A_\mu^1+g_2q_2A_\mu^2)^2.
  \]
  With \(N=\sqrt{g_1^2q_1^2+g_2^2q_2^2}\), define
  \[
    A_\mu^M=\frac{g_1q_1A_\mu^1+g_2q_2A_\mu^2}{N},
    \qquad
    A_\mu^0=\frac{-g_2q_2A_\mu^1+g_1q_1A_\mu^2}{N}.
  \]
  Then \(m_M^2=v^2N^2\) in the usual complex-scalar normalization,
  whereas \(A^0\) is massless.  The unbroken transformation satisfies
  \(q_1\alpha_1+q_2\alpha_2=0\); its generator is proportional to
  \(q_2Q_1-q_1Q_2\).

  \SupplementarySolution{10}{Fermion Mass and Higgs Coupling}
  Substitution of \(H=(v+h)/\sqrt2\) gives
  \[
    \mathcal L_Y
      =-\frac{yv}{\sqrt2}\bar\psi\psi
       -\frac{y}{\sqrt2}h\bar\psi\psi.
  \]
  Hence
  \[
    \boxed{m_\psi=\frac{yv}{\sqrt2},\qquad
      \mathcal L_{h\bar\psi\psi}
       =-\frac{m_\psi}{v}h\bar\psi\psi.}
  \]
  The Higgs coupling is therefore not a new independent parameter once
  the fermion mass and \(v\) are known.

  \SupplementarySolution{11}{Electric Charge from \(T_3+Y\)}
  Since \(Q(\nu_L,e_L)=(0,-1)\) and
  \(T_3=(+\tfrac12,-\tfrac12)\),
  \[
    Y_L=-\frac12.
  \]
  The singlet has \(T_3=0\) and \(Q=-1\), so \(Y_{e_R}=-1\).
  Requiring \(Q(H^+,H^0)=(1,0)\) gives \(Y_H=+\tfrac12\).
  The hypercharge of
  \(\bar LHe_R\) is
  \[
    -Y_L+Y_H+Y_{e_R}
      =\frac12+\frac12-1=0.
  \]
  The doublet indices contract to an \(SU(2)\) singlet, so the Yukawa
  interaction is gauge invariant.

  \SupplementarySolution{12}{The Dimension-Five Neutrino Operator}
  A convenient form, with flavor labels suppressed, is
  \[
    \boxed{\mathcal L_5
      =\frac{c}{\Lambda}
       \bigl(L^{\mathsf T}C\,i\sigma_2H\bigr)
       \bigl(H^{\mathsf T}i\sigma_2L\bigr)
       +\mathrm{H.c.}}
  \]
  with the spinor contraction understood.  Each \(LH\) factor is an
  \(SU(2)\) singlet and has zero hypercharge.  Since
  \([L]=3/2\) and \([H]=1\), the operator has dimension five.
  With \(\langle H\rangle=(0,v/\sqrt2)^{\mathsf T}\), it becomes a
  Majorana mass term for \(\nu_L\), with
  \[
    \boxed{m_\nu=\frac{|c|v^2}{2\Lambda}}
  \]
  up to a convention-dependent factor absorbed into the symmetric
  flavor coefficient.

  \SupplementarySolution{13}{Neutral-Current Couplings}
  Write
  \[
    Z_\mu=\cos\theta_W W_\mu^3-\sin\theta_WB_\mu,
    \qquad
    g\sin\theta_W=g'\cos\theta_W=e.
  \]
  Substituting \(Y=Q-T_3\) into \(gT_3W^3+g'YB\) gives the \(Z\)
  generator
  \[
    \frac{g}{\cos\theta_W}(T_3-Q\sin^2\theta_W).
  \]
  Separating left and right chiralities yields
  \[
    g_V^f=T_3^f-2Q_fs_W^2,\qquad g_A^f=T_3^f.
  \]
  Thus
  \[
    (g_V^\ell,g_A^\ell)
      =(-\tfrac12+2s_W^2,-\tfrac12),
    \qquad
    (g_V^\nu,g_A^\nu)=(\tfrac12,\tfrac12).
  \]

  \SupplementarySolution{14}{The Tree-Level Mass Relation}
  The vacuum derivative is
  \[
    D_\mu\langle H\rangle
      =-\frac{i v}{2\sqrt2}
       \begin{pmatrix}
         g(W_\mu^1-iW_\mu^2)\\
         -gW_\mu^3+g'B_\mu
       \end{pmatrix}.
  \]
  Hence
  \[
    m_W^2=\frac{g^2v^2}{4},
    \qquad
    \mathcal L_{\rm neutral,mass}
      =\frac{v^2}{8}(gW^3-g'B)^2.
  \]
  Diagonalizing the neutral matrix gives
  \[
    Z=\cos\theta_WW^3-\sin\theta_WB,\qquad
    A=\sin\theta_WW^3+\cos\theta_WB,
  \]
  with \(\tan\theta_W=g'/g\).  Therefore
  \[
    \boxed{m_Z=\frac v2\sqrt{g^2+g'^2},\quad
      m_A=0,\quad m_W=m_Z\cos\theta_W.}
  \]

  \SupplementarySolution{15}{Determining the Electroweak Scale}
  Each charged-current vertex contributes \(g/\sqrt2\) times a chiral
  projector.  At low momentum the \(W\) propagator reduces to
  \(1/m_W^2\), and comparison with the conventional four-fermion
  interaction gives
  \[
    \frac{G_F}{\sqrt2}=\frac{g^2}{8m_W^2}.
  \]
  Using \(m_W=gv/2\),
  \[
    \frac{G_F}{\sqrt2}=\frac1{2v^2},
    \qquad
    \boxed{v=(\sqrt2G_F)^{-1/2}.}
  \]
  Numerically,
  \[
    \boxed{v\simeq246.22\ {\rm GeV}.}
  \]

  \SupplementarySolution{16}{A Leptonic \(Z\) Width}
  Squaring the neutral-current amplitude, summing final spins and colors,
  and using the two-body massless phase space gives
  \[
    \Gamma(Z\to f\bar f)
      =N_c\frac{g^2m_Z}{48\pi\cos^2\theta_W}
       [(g_V^f)^2+(g_A^f)^2].
  \]
  The relations
  \(g^2/\cos^2\theta_W=4m_Z^2/v^2\) and
  \(v^{-2}=\sqrt2G_F\) yield
  \[
    \boxed{\Gamma
      =N_c\frac{G_Fm_Z^3}{6\pi\sqrt2}
       [(g_V^f)^2+(g_A^f)^2].}
  \]
  For a charged lepton,
  \(g_A=-1/2\) while \(g_V=-1/2+2s_W^2\) is close to zero.
  The interaction is therefore dominantly axial and strongly violates
  parity.

  \SupplementarySolution{17}{A Heavy Doublet and the \(T\) Parameter}
  Put \(m_1=m+\delta m/2\) and \(m_2=m-\delta m/2\).  At
  \(\delta m=0\), the logarithmic term has the limiting value
  \(2m^2\), which cancels \(m_1^2+m_2^2=2m^2\).  Expanding the logarithm
  and the rational prefactor gives
  \[
    \boxed{\Delta m^2
      =\frac43(\delta m)^2
       +O\!\left(\frac{(\delta m)^4}{m^2}\right).}
  \]
  Thus an exactly degenerate weak doublet respects custodial symmetry
  and does not contribute to \(T\), whereas the leading violation is
  positive and quadratic in the mass splitting, multiplied by \(N_c\)
  in the full self-energy correction.

  \SupplementarySolution{18}{The One-Loop Electroweak Scale of \(a_\ell\)}
  A magnetic dipole operator flips chirality, so the amplitude supplies
  one lepton mass.  Its conventional normalization contains another,
  while the heavy weak propagator supplies \(1/m_W^2\).  Consequently
  \(a_\ell^{\rm EW}\propto G_Fm_\ell^2\).

  The bracket is flavor-independent at this accuracy, so
  \[
    \frac{a_\mu^{\rm EW}}{a_e^{\rm EW}}
      =\left(\frac{m_\mu}{m_e}\right)^2
      \simeq(206.77)^2
      \simeq\boxed{4.28\times10^4}.
  \]
  This mass enhancement is why weak and hadronic effects are far more
  visible in the muon anomaly than in the electron anomaly.

  \SupplementarySolution{19}{One Standard-Model Family in \(SU(5)\)}
  The decompositions are
  \[
    \boxed{\overline{\mathbf5}
      =(\overline{\mathbf3},\mathbf1,+\tfrac13)
       \oplus(\mathbf1,\mathbf2,-\tfrac12)}
  \]
  and
  \[
    \boxed{\mathbf{10}
      =(\mathbf3,\mathbf2,+\tfrac16)
       \oplus(\overline{\mathbf3},\mathbf1,-\tfrac23)
       \oplus(\mathbf1,\mathbf1,+1).}
  \]
  The first contains \(d^c\) and \(L=(\nu,e)_L\).  The antisymmetric
  \(\mathbf{10}\) contains \(Q=(u,d)_L\), \(u^c\), and \(e^c\).
  The dimensions \(5+10=15\) equal the number of left-handed Weyl fields
  in one family without a right-handed neutrino.

  \SupplementarySolution{20}{The \(5/3\) Hypercharge Normalization}
  In the fundamental,
  \[
    \operatorname{Tr}Y^2
      =3\left(\frac13\right)^2
       +2\left(\frac12\right)^2
      =\frac56.
  \]
  A conventionally normalized \(SU(5)\) generator has trace square
  \(1/2\), so
  \[
    T_Y=\sqrt{\frac35}\,Y.
  \]
  Equality of the covariant derivative,
  \(g_1T_Y=g'Y\), then requires
  \[
    \boxed{g_1=\sqrt{\frac53}\,g',
      \qquad \alpha_1=\frac53\alpha_Y.}
  \]
  The factor is a generator normalization, not an additional physical
  interaction.

  \SupplementarySolution{21}{Solving One-Loop Unification}
  Set \(A_i=\alpha_i^{-1}(\mu)\) and
  \(L=\ln(M_U/\mu)\).  Subtraction eliminates the unified coupling:
  \[
    A_i-A_j=\frac{b_i-b_j}{2\pi}L.
  \]
  Therefore
  \[
    \boxed{L=\frac{2\pi(A_i-A_j)}{b_i-b_j},\qquad
      \alpha_U^{-1}
       =\frac{b_iA_j-b_jA_i}{b_i-b_j}.}
  \]
  The third coupling must satisfy
  \(A_k=\alpha_U^{-1}+b_kL/(2\pi)\).  Failure of this equality measures
  the need for thresholds, additional matter, higher-loop corrections,
  or the failure of the proposed unification model.

  \SupplementarySolution{22}{Why Complete Multiplets Preserve a Meeting}
  Only coefficient differences enter
  \[
    \alpha_i^{-1}(\mu)-\alpha_j^{-1}(\mu)
      =\frac{b_i-b_j}{2\pi}\ln\frac{M_U}{\mu}.
  \]
  If every \(b_i\) receives the same shift \(\Delta b\), these
  differences and the extracted \(M_U\) are unchanged.  Each individual
  equation does change:
  \[
    \alpha_U^{-1}\longrightarrow
    \alpha_U^{-1}
      -\frac{\Delta b}{2\pi}\ln\frac{M_U}{\mu}
  \]
  for fixed low-energy couplings in this sign convention.
  A degenerate complete representation of the unified group has the same
  total index in every subgroup after proper normalization, and hence
  supplies precisely such a common shift.  Split or incomplete
  multiplets need not.

  \SupplementarySolution{23}{Dimension-Six Proton Decay}
  Tree-level \(X\) exchange gives
  \[
    \mathcal L_{\rm eff}
      \sim\frac{g_U^2}{M_X^2}(qq)(q\ell).
  \]
  The coefficient has dimension \(-2\).  A proton-decay amplitude is
  therefore of order \(g_U^2m_p^2/M_X^2\), apart from dimensionless
  group factors and a hadronic matrix element.  Multiplying its square
  by the one-body decay scale \(m_p\) gives
  \[
    \boxed{\Gamma_p\sim
      \frac{\alpha_U^2m_p^5}{M_X^4},\qquad
      \tau_p\sim\frac{M_X^4}{\alpha_U^2m_p^5}.}
  \]
  The fourth power makes nucleon lifetime predictions exceptionally
  sensitive to the heavy gauge-boson mass.

  \SupplementarySolution{24}{Threshold Corrections}
  With thresholds,
  \[
    A_i-A_j
      =\frac{b_i-b_j}{2\pi}L+(\delta_i-\delta_j).
  \]
  Thus
  \[
    L=\frac{2\pi[(A_i-A_j)-(\delta_i-\delta_j)]}{b_i-b_j}.
  \]
  Nonuniversal differences \(\delta_i-\delta_j\) move the inferred
  meeting scale and can change whether the third coupling meets the
  first two.  Once \(L\) is fixed, any one equation gives
  \(\alpha_U^{-1}=A_i-b_iL/(2\pi)-\delta_i\); a common part of all
  \(\delta_i\) therefore shifts only the unified coupling.

  For a degenerate complete multiplet,
  \(b_1^{(r)}=b_2^{(r)}=b_3^{(r)}\) after summing its components, and all
  masses give the same logarithm.  Its threshold is common and cannot
  alter coupling differences, so it cannot repair a mismatch.

  \SupplementarySolution{25}{London Screening}
  Varying the matter part of the free energy with respect to
  \(\mathbf A\) gives
  \[
    \mathbf J
      =\rho_s\frac q\hbar
       \left(\boldsymbol\nabla\theta-\frac q\hbar\mathbf A\right).
  \]
  In a simply connected static region choose a gauge with constant
  \(\theta\).  Then
  \[
    \mathbf J=-\frac{\rho_sq^2}{\hbar^2}\mathbf A.
  \]
  Taking the curl of \(\boldsymbol\nabla\times\mathbf B=\mathbf J\),
  using \(\boldsymbol\nabla\cdot\mathbf B=0\), gives
  \[
    (\nabla^2-\lambda^{-2})\mathbf B=0,\qquad
    \boxed{\lambda=\frac{\hbar}{|q|\sqrt{\rho_s}}}
  \]
  in the units of the exercise.  A surface field therefore decays
  exponentially over the penetration depth \(\lambda\).

  \SupplementarySolution{26}{Flux Quantization for Charge \(q\)}
  Vanishing bulk current requires
  \[
    \boldsymbol\nabla\theta=\frac q\hbar\mathbf A.
  \]
  Single-valuedness gives
  \(\oint\nabla\theta\cdot d\mathbf x=2\pi n\), so
  \[
    \boxed{\Phi=\oint\mathbf A\cdot d\mathbf x
      =\frac{2\pi\hbar}{q}n=\frac hq n.}
  \]
  The sign records orientation.  For a Cooper-pair condensate
  \(q=-2e\), the magnitude of the smallest nonzero flux is
  \[
    \boxed{\Phi_0=\frac h{2e}=\frac{\pi\hbar}{e}.}
  \]

  \SupplementarySolution{27}{The ac Josephson Relation}
  Differentiating the junction energy gives the charge current
  \[
    I=\frac{qE_J}{\hbar}\sin\delta
  \]
  up to the orientation convention.  Gauge invariance and the phase
  equation of motion give
  \[
    \dot\delta=-\frac q\hbar\Delta V.
  \]
  Thus a constant voltage produces a sinusoidal current of angular
  frequency and ordinary frequency
  \[
    \boxed{\omega_J=\frac{|q\Delta V|}{\hbar},\qquad
      \nu_J=\frac{|q\Delta V|}{2\pi\hbar}
       =\frac{|q\Delta V|}{h}.}
  \]
  For \(q=-2e\), this is the standard \(2e|\Delta V|/h\).

  \SupplementarySolution{28}{A Phase Slip Changes the Fluxoid}
  On any contour where the order parameter never vanishes, its
  normalized phase defines a continuous map \(S^1\to S^1\).  The
  integer
  \[
    n=\frac1{2\pi}\oint\nabla\theta\cdot d\mathbf x
  \]
  is a homotopy invariant and cannot change under a continuous
  deformation.  During a phase slip the modulus must therefore vanish
  at some spacetime point, making the phase undefined and allowing
  \(n\to n\pm1\).  Since \(\Phi=nh/q\) deep in the ring, one slip changes
  the flux by
  \[
    \boxed{\Delta\Phi=\pm\frac hq}
  \]
  with sign set by the direction of the slip.

  \SupplementarySolution{29}{Several Commensurate Condensates}
  Write a gauge transformation as
  \(\theta_i\to\theta_i+n_i\alpha\), with
  \(\alpha\simeq\alpha+2\pi\).  It leaves every condensate invariant when
  \(n_i\alpha\in2\pi\mathbb Z\) for all \(i\).  If
  \(d=\gcd(n_i)\), the solutions are
  \(\alpha=2\pi k/d\), so the residual group is
  \[
    \boxed{\mathbb Z_d.}
  \]
  The simultaneous flux conditions are
  \(n_iq_0\Phi/\hbar\in2\pi\mathbb Z\).  Their smallest common nonzero
  solution is
  \[
    \boxed{\Phi_0=\frac h{d q_0}}
  \]
  in magnitude.  Josephson harmonics have frequencies
  \(|n_iq_0\Delta V|/h\); all integer combinations share the fundamental
  frequency
  \[
    \boxed{\nu_0=\frac{d|q_0\Delta V|}{h}.}
  \]

  \SupplementarySolution{30}{Incommensurate Charges and Relative Phases}
  Each condensate separately requires
  \[
    \Phi\in\frac h{q_1}\mathbb Z,
    \qquad
    \Phi\in\frac h{q_2}\mathbb Z.
  \]
  If \(q_1/q_2\) is irrational, equality
  \(n_1/q_1=n_2/q_2\) has no nonzero integer solution.  Therefore
  \[
    \boxed{\frac h{q_1}\mathbb Z\cap
      \frac h{q_2}\mathbb Z=\{0\}.}
  \]
  Under a gauge transformation,
  \(\theta_i\to\theta_i+q_i\Lambda/\hbar\), so
  \[
    \boxed{\Theta_{\rm rel}=q_2\theta_1-q_1\theta_2}
  \]
  is gauge invariant.  Without a locking interaction it is an
  independent collective coordinate.  Nonzero common flux sectors are
  absent, while a junction current contains oscillations at
  \(|q_1\Delta V|/h\) and \(|q_2\Delta V|/h\).  Their irrational ratio
  makes the total current quasiperiodic rather than periodic with one
  Josephson frequency.
\end{enumerate}
